LLM-based two-agent solver-checker systems communicate prolifically---a single math problem can generate dozens of messages---yet which are structurally necessary remains open. Existing approaches operate at different granularities: token-level compression~\cite{jiang_llmlingua_2023, chevalier_autocompressors_2023}, runtime trajectory filtering~\cite{lin_supervisoragent_2026}, agent-level pruning~\cite{wang_agentdropout_2025, zhang_agentprune_2025}, message-level suppression~\cite{li_agp_2025, zeng_s2mad_2025}, and system-level optimization~\cite{chen_optima_2025}. Yet none perform \emph{functional analysis}---understanding what communicative role each message plays---nor validate whether functional classification predicts which messages can be safely removed.

An intuitive hypothesis is common: high-information messages are important; low-information acknowledgments are expendable. This conflates \emph{perturbation sensitivity}---how much removing a message disrupts flow---with \emph{informational importance}---whether it carries task-relevant content. Pruning methods that score by content richness~\cite{wang_agentdropout_2025, zeng_s2mad_2025} implicitly assume these correlate. Yet conversation analysis shows backchannels~\cite{yngve_getting_1970, schegloff_discourse_1982, jefferson_acknowledgment_1984} are structurally load-bearing despite negligible propositional content~\cite{clark_contributing_1989}. Whether LLM agents exhibit analogous dependencies has never been tested.

\paragraph{The gap.} Three layers of evidence point in contradictory directions: (1)~\emph{information content} suggests acknowledgments are dispensable; (2)~\emph{individual perturbation} reveals they are structurally critical (91\% high disruption); (3)~\emph{aggregate pruning} shows they are safely removable ($\Delta \approx 0$\,pp). No prior work has tested this three-way paradox in LLM agents nor explained why individually critical messages become safely removable in aggregate~\cite{rizvi_martel_communication_2026}.

\paragraph{Contributions.} To our knowledge, we present the first perturbation-based empirical study of message-level pruning in two-agent solver-checker LLM systems, spanning two primary backbones (Qwen, Llama; Gemma exploratory), \numValidatedDomains{} validated task domains (with ARC and MATH as supplementary), and three topologies (15 configurations at 7B, with 14B and preliminary 72B scale probes):

\begin{enumerate}[nosep, leftmargin=*]
    \item \textbf{Perturbation-sensitivity dissociation} (\S\ref{sec:impact_gradient}--\ref{sec:three_modes}). Messages most disruptive to perturb individually can be pruned in aggregate without significant accuracy loss. We conjecture that self-reinforcing ack cascades create structural redundancy that dissolves under aggregate removal (Qwen 90.7\%, Llama 61.9\%)---a mechanism distinct from human backchannel dynamics, requiring causal verification. This finding is enabled by a reusable analysis protocol comprising a 5-class taxonomy, structured annotation, and perturbation-based impact measurement.

    \item \textbf{Analysis protocol} (\S\ref{sec:taxonomy_method}--\ref{sec:statistical_framework}). A 5-class communicative function taxonomy validated by cross-model agreement ($\kappa{=}0.721$) and human annotation ($\kappa{=}0.76$), a structured annotation protocol (constrained JSON, cross-model validation), and a perturbation-based impact measurement pipeline with non-inferiority and TOST equivalence testing---designed to be applicable to any two-agent LLM system for message-level communication analysis. The taxonomy serves a dual role: as a scientific lens for understanding communication structure, and as a label generator enabling automated pruning without per-message annotation at deployment.

    \item \textbf{Pruning: from oracle to deployment} (\S\ref{sec:e2e_pruning}, \ref{sec:boundary}). Oracle and annotation-free heuristic pruning reduce inference calls by 36--56\% with non-inferiority at $\delta = -3$\,pp ($p < 0.001$) and supplementary Qwen-specific TOST equivalence at $\pm 1$\,pp (oracle $p = 0.048$, heuristic $p = 0.007$; 10 seeds each), validated across two backbone families and \numValidatedDomains{} task domains under greedy decoding. Boundary conditions---heuristic non-transferability (Qwen3), domain mismatch (MATH), and backbone-dependent stochastic decoding sensitivity---motivate empirical screening ($N \geq 200$) and a learned-classifier path beyond hand-crafted rules.
\end{enumerate}

\paragraph{Scope.} Claims are validated under greedy decoding; stochastic decoding ($t{=}0.7$) is validated for Qwen (10 seeds, TOST $p < 0.001$ at $\delta = -3$\,pp) with Llama showing a marginal degradation trend (\S\ref{sec:results}, \S\ref{sec:discussion}). Our message-level analysis complements token-level compression~\cite{jiang_llmlingua_2023} and agent-level pruning~\cite{wang_agentdropout_2025}, bridging theoretical bounds~\cite{rizvi_martel_communication_2026} and practical pruning. The taxonomy serves as an analytical lens; heuristic detectors that make it unnecessary at deployment are products of the analysis (\S\ref{sec:related_work}).
